% !TEX root = ../bb_AiRa_Kappaleet.tex

% ö = \%c3\%b6
% ä = \%C3\%A4

\xdef\sectiontitle{\IIIsectiontitle} %aiheuttama
\TOCrefs

%%%%%%%%%%%%%%%
\begin{frame}[label=3_7_a]%[label=3_7_TOC1]
\frametitle{\texorpdfstring{\culine[\sectiontitleulinecolor]{\sectiontitle}}{\sectiontitle} }
\label{\TOCname}

\vspace{-0.5cm}\label{\TOCname} % vertical spacing is off, 0.7cm doesn't work

\begin{itemize}[leftmargin=0.5cm,itemsep=\chapterTOCitemsep]
    \item[3.1]<1-> \hyperlink{3_1_a}{\IIIaSubsectiontitle}
    \item[3.2]<1-> \hyperlink{3_2_a}{\IIIbSubsectiontitle}	
    \item[3.3]<1-> \hyperlink{3_3_a}{\IIIcSubsectiontitle}	
    \item[3.4]<1-> \hyperlink{3_4_a}{\IIIdSubsectiontitle}
    \item[3.5]<1-> \hyperlink{3_5_a}{\IIIeSubsectiontitle}
    \item[3.6]<1-> \hyperlink{3_6_a}{\IIIfSubsectiontitle}
    \item[3.7]<1-> \hyperlink{3_7_a}{\CDAlert[blue]{\IIIgSubsectiontitle}}
\end{itemize}

%\endofslide 
\end{frame}
%%%%%%%%%%%%%%%

\xdef\sectiontitle{\IIIgSubsectiontitle} 


%%%%%%%%%%%%%%%
\begin{frame}
\frametitle{\texorpdfstring{\culine[\sectiontitleulinecolor]{\sectiontitle}}{\sectiontitle} }
\framesubtitle{NMR}
\appear{}

\begin{itemize}
	\item[] \CDAlert[fuchsia]{Ydinmagneettinen resonanssi (NMR)}:
	
	\item Kuten \CDAlert[red]{Zeemanin ilmiö} atomin elektroneilla\appear{,  \Alert[red]{myös ytimen energiatasot jakaantuvat ulkoisessa magneettikentässä}}

\item Ytimen \CDAlert[blue]{magneettinen momentti} on  
\[
\vec{\mu}  \equal \appear{\gamma} \appear{\vec{I}} \appear{\,, \qquad \mu_{N}}  \equal \frac{e \hbar}{2 m_{p}} \appear{\,, \qquad \text{missä}} \qquad \appear{|\, \vec{I} \,|}  \equal  \appear{I(I+1)^{1 / 2} \hbar}
\]

\item Ulkoisessa magneettikentässä olevan magneettimomentin \CDAlert[fuchsia]{potentiaalienergia} on
\[
U  \equal \appear{-\vec{\mu}} \appear{\cdot \vec{B}}
\]

\item \CDAlert[fuchsia]{Energiaero} jakaantuneiden tilojen välillä on
\[
\Delta E  \equal \appear{2} \appear{( \mu_{z} )_{\text {protoni}}} \appear{B}
\]
\end{itemize}

\endofslide 
%\endofsection
\end{frame}
%%%%%%%%%%%%%%%

%%%%%%%%%%%%%%%
\begin{frame}
\frametitle{\texorpdfstring{\culine[\sectiontitleulinecolor]{\sectiontitle}}{\sectiontitle} }
\framesubtitle{NMR}
\appear{}

\seminormal 
\begin{itemize}
	\item Kun vetyatomeja \appear{(esim.\ veden protoneja)} \appear{\Alert[blue]{viritetään sähkömagneettisella säteilyllä}}\appear{, eli fotoneilla $h f=\Delta E$}\appear{, tapahtuu \CDAlert[blue]{resonanssiabsorptiota}, ja \CDAlert[blue]{vastaavaa emissiota}}

\item $1 \mathrm{\;T}$ magneettikentässä\appear{, energiaero $\Delta E=1,76 \times 10^{-7} \mathrm{\;eV}$} \appear{vastaa resonanssitaajuutta $f=42,5$ $\mathrm{\;MHz}$}  \appear{(radiotaajuus)}

\end{itemize}

% 6.5 cm = midway, 10 cm = 3/4 
\vspace{2.5mm}
\begin{minipage}{7.5cm}
\begin{itemize}
	\item NMR:ää sovelletaan \appear{materiaalien \Alert[fuchsia]{sisäisen \CDAlert[fuchsia]{magneettisen rakenteen}} tutkimiseen}\appear{, ja erityisesti elävien kudosten \Alert[fuchsia]{magneettiresonanssikuvaukseen (\CDAlert[fuchsia]{MRI})}}\appear{, minimoiden kudoksien lämpenemisen}
\appear{\xdef\loc{south east}
\begin{tikzpicture}[remember picture,overlay]% anchor=south
    \node (A) [yshift=-0.35*\figYshift,xshift=\figXshift, anchor=\loc] at (current page.\loc) {\includegraphics[width=5.8cm]{Figures_booklet/SOM_12_Nuclear_physics_figures/SOM_Nuclear_physics_Figure_17.png}};% 8cm max height
\end{tikzpicture}\CR{}}
\item NMR on parhaiten tunnettu MR-kuvauksesta\appear{, mutta sitä käytetään paljon myös tutkimuksessa}\appear{, kuten \CDAlert[fuchsia]{mo-} \CDAlert[fuchsia]{nimutkaisten kemiallisten molekyylien} \Alert[fuchsia]{rakenteiden karakterisointiin} }
\end{itemize}
\end{minipage}
%\vspace{2.5mm}

\endofslide 
%\endofsection
\end{frame}
%%%%%%%%%%%%%%%

%%%%%%%%%%%%%%%
\begin{frame}
\frametitle{\texorpdfstring{\culine[\sectiontitleulinecolor]{\sectiontitle}}{\sectiontitle} }
\framesubtitle{NAA}
\appear{}

% Neutron Activation Analysis (NAA)
\begin{itemize}[itemsep=2.2mm]
	\item[] \CDAlert[fuchsia]{Neutroniaktivaatioanalyysi (NAA)}:
	\item ${ }_{Z}^{A} M$ sisältävää näytettä pommitetaan hitailla neutroneilla 

\item Nuklidi ${ }^{A+1}_{~\;Z}M$ \appear{voi muodostua reaktiosta} $\appear{{ }_{Z}^{A} M} \appear{(}\appear{n,} \appear{\gamma)}\appear{^{A+1} M}$\appear{, mutta se \CDAlert[blue]{on radioaktiivinen}}\appear{, eli se emittoi beetahiukkasia tai gammasäteitä}

\item \CDAlert[blue]{Muodostuvaa säteilyä mittaamalla}\appear{, ${ }^{A+1}_{~\;Z} M$ voidaan tunnistaa}\appear{, ja lisäksi isotoopin ${ }_{Z}^{A} M $ pitoisuus näytteessä voidaan määrittää}

\item Tämä mahdollistaa näppärästi \CDAlert[blue]{kemiallisen analyysin}\appear{, jopa erilaisia isotooppeja voidaan erottaa toisistaan}\appear{, hajottamatta näytettä}

\item \CDAlert[fuchsia]{Neutroniaktivaatioanalyysin (NAA)} avulla voidaan analysoida pieniä määriä ainetta\appear{, ja sitä käytetään usein kun 'muut fysikokemiallisen menetelmät eivät ole tarpeeksi herkkiä'}

\item Menetelmän heikkoutena on radioaktiivisten näytteiden \CDAlert[red]{hävittäminen}

\end{itemize}

\endofslide 
%\endofsection
\end{frame}
%%%%%%%%%%%%%%%

%%%%%%%%%%%%%%%
\begin{frame}
\frametitle{\texorpdfstring{\culine[\sectiontitleulinecolor]{\sectiontitle}}{\sectiontitle} }
\framesubtitle{NAA}
\appear{}

\seminormal
\begin{itemize}
	\item Kun nuklideja on säteilytetty termisillä neutroneilla, niiden aktiivisuus noudattaa yhtälöä
\[
R(t)  \equal \appear{\lambda} \appear{N(t)}  \equal \appear{R_{0}} \appear{(1-\rme^{-\lambda t} )}
\]
\appear{missä $\lambda=$ hajoamisvakio} \appear{joka riippuu ytimestä ${ }^{A+1}_{~\;Z}M$}\appear{, ja $R_{0}$ on isotoopin  muodostumisnopeus} \appear{($g$ aiemmissa luennoissa)}

\item Yhtälössä\appear{ $R(t)$ mitataan ja $R_{0}$ voidaan periaatteessa laskea tunnettujen intensiteettien $I$ avulla} \appear{(riippuu mittausjärjestelystä)} \appear{ja vuorovaikutusalasta} \appear{(riippuu reaktiosta joka johtaa nuklidiin $ { }^{A+1}_{~\;Z}M$).} \appear{$I$ on neutronivuo (esim.\ yksiköissä neutronia/s$\cdot\mathrm{m}^{2}$)}

\item Yleisesti ottaen ytimen ${ }^{A+1}_{~\;Z}M$ \CDAlert[blue]{törmäys/tuotantonopeus} on:
\[
R_{0}  \equal \appear{N_{0} \sigma I} \qquad \appear{, \qquad [\sigma I]}  \equal \appear{m^{2}} \frac{\text {neutronia}}{s \cdot m^{2}}  \equal \frac{\text {neutronia}}{s}
\]\vspace{2.0mm}
\appear{missä $N_{0}$ on nuklidien ${ }_{Z}^{A} M$ lukumäärä näytteessä}
\end{itemize}

\endofslide 
%\endofsection
\end{frame}
%%%%%%%%%%%%%%%

%%%%%%%%%%%%%%%
\begin{frame}
\frametitle{\texorpdfstring{\culine[\sectiontitleulinecolor]{\sectiontitle}}{\sectiontitle} }
\framesubtitle{NAA}
\appear{}

\begin{itemize}
	\item Näin saadaan $\appear{R(t)}  \equal \appear{\lambda N(t)}  \equal \appear{N_{0} \sigma I (1-\rme^{-\lambda t} )}$ \appear{ja jos nuklidin puoliintumisaika on lyhyt}\appear{, aktiivisuus saturoituu $\appear{R(t)}\appear{=R(t=\infty)}\appear{=N_{0} \sigma I}$}

\item Taulukko näyttää \CDAlert[red]{saturaatioaktiivisuudet $R(\infty)$} \appear{(per näytemilligramma)} \appear{muutamille isotoopeille}
\appear{\xdef\loc{south east}%
\begin{tikzpicture}[remember picture,overlay]% anchor=south
    \node (A) [yshift=-0.25*\figYshift,xshift=\figXshift, anchor=\loc] at (current page.\loc) {\small \def\arraystretch{1.4}$
\begin{array}{c|c|c}
 & & R(\infty)  \\
{ }_{Z}^{A} \mathrm{M} & { }^{A+1}_{~\;Z} \mathrm{M} & \text {hajoamista} / \mathrm{min} \cdot \mathrm{mg} \\
\hline
{ }^{55} \mathrm{Mn} & { }^{56} \mathrm{Mn} & 8,8 \times 10^{6} \\
{ }^{63} \mathrm{Cu} & { }^{64} \mathrm{Cu} & 1,7 \times 10^{6} \\
{ }^{127} \mathrm{I} & { }^{128} \mathrm{I} & 1,6 \times 10^{6} \\
{ }^{197} \mathrm{Au} & { }^{198} \mathrm{Au} & 1,7 \times 10^{7}
\end{array}
$};% 8cm max height
\end{tikzpicture}}
\item Atomien ${ }_{Z}^{A} M$ lukumäärä näytteessä on
\[
N_{0}  \equal \fraca{R(t=\infty)}{\sigma I} \appear{\,,}
\]
\appear{ja näytteessä olevan ${ }_{Z}^{A} M$:n massa on}
\end{itemize}

% 6.5 cm = midway, 10 cm = 3/4 
\vspace{2.5mm}
\begin{minipage}{7.8cm}
\begin{itemize}
	\item<.->[] 
	\[
 \appear{m ({ }_{Z}^{A} M)}  \equal \frac{N_{0}}{N_{A}} \appear{M({ }_{Z}^{A} M)}  \equal \frac{R(t=\infty)}{N_{A} \sigma I} \appear{M({ }_{Z}^{A} M) \,,}
\]
\item[] missä $M({ }_{Z}^{A} M)$ on alkuaineen atomipaino \appear{ja $N_{A}$ on Avogadron luku}	
\end{itemize}
\end{minipage}
%\vspace{2.5mm}

\endofslide 
%\endofsection
\end{frame}
%%%%%%%%%%%%%%%

%%%%%%%%%%%%%%%
\begin{frame}
\frametitle{\texorpdfstring{\culine[\sectiontitleulinecolor]{\sectiontitle}}{\sectiontitle} }
\framesubtitle{Ajoitus}
\appear{}

\begin{itemize}
	\item[] \CDAlert[fuchsia]{Radioaktiivinen ajoitus}:
\end{itemize}
	
%Radiaoctive Dating
\semismall
\begin{itemize}
	%\item[] \CDAlert[fuchsia]{Radiaoctive Dating}:
	\item Luonnossa on radioaktiivisia isotooppeja:
\begin{itemize}
	\item[(1)] Kolme alfaaktiivista alkuainetta ketjuineen \appear{(\CDAlert[blue]{torium}}\appear{, \CDAlert[green]{uraani}}\appear{, \CDAlert[red]{aktinium})}
\item[(2)] Yksittäisiä pitkäikäisiä isotooppeja\appear{, kuten ${ }^{40} \mathrm{K}$} \appear{($T_{1 / 2}=1,25 \times 10^{9} \mathrm{\;a}$)}
\item[(3)] Yksittäisiä radioaktiivisia isotooppeja muodostuu \CDAlert[fuchsia]{kosmisten säteiden} \CDAlert[fuchsia]{pommituksessa}
\end{itemize}

\item Geologiseen ajoitukseen käytetään \appear{${ }^{40} \mathrm{K}$} \appear{ja ${ }^{232} \mathrm{Th} ~(T_{1 / 2}=1,24 \times 10^{10} \mathrm{\;a})$} \appear{(\CDAlert[blue]{"vanhoihin" ~kiviin}}\appear{, tai "uudempiin" ~kiviin)}

\item Kahden isotoopin \CDAlert[fuchsia]{nykyinen pitoisuussuhde} mitataan\appear{, joista toinen on radioaktiivinen tai radioaktiivisen hajoa-}%time dependent)
\end{itemize}

% 6.5 cm = midway, 10 cm = 3/4 
\vspace{-0.3mm}
\begin{minipage}{6.45cm}
\begin{itemize}
	\item[] misen stabiili lopputuote, (jolloin pi\-toi\-suus\-suhde on ajastariippuva)\appear{, ja verrataan \CDAlert[red]{aiempaan} \CDAlert[red]{pitoisuussuhteeseen}}\appear{, jonka on tiedetty olevan voimassa kun aine muodostui}
\end{itemize}
\end{minipage}
%\vspace{2.5mm}


\xdef\loc{south east}
\begin{tikzpicture}[remember picture,overlay] % anchor=south
    \node (A) [yshift=-0.25*\figYshift,xshift=\figXshift, anchor=\loc] at (current page.\loc) {\footnotesize \def\arraystretch{1.15}$
\begin{array}{c|c|c|c}
\text {Nuklidi} & t_{1 / 2}(a) & \text {Pitoisuus (\%) } & \text {Tytär} \\
\hline
{ }^{14} \mathrm{C} & 5730 & 1.35 \times 10^{-10} & { }^{14} \mathrm{N} \\
{ }^{40} \mathrm{K} & 1,25 \times 10^{9} & 0,0117 & { }^{40} \mathrm{A} \\
{ }^{87} \mathrm{Rb} & 4,88 \times 10^{10} & 27,83 & { }^{87} \mathrm{Sr} \\
{ }^{147} \mathrm{Sm} & 1,06 \times 10^{11} & 15,0 & { }^{143} \mathrm{Nd} \\
{ }^{176} \mathrm{Lu} & 3,59 \times 10^{10} & 2,59 & { }^{176} \mathrm{Hf} \\
{ }^{187} \mathrm{Re} & 4,30 \times 10^{10} & 62,60 & { }^{18} \mathrm{Os}
\end{array}
$}; % 8cm max height
\end{tikzpicture}

\endofslide 
%\endofsection
\end{frame}
%%%%%%%%%%%%%%%

%%%%%%%%%%%%%%%
\begin{frame}
\frametitle{\texorpdfstring{\culine[\sectiontitleulinecolor]{\sectiontitle}}{\sectiontitle} }
\framesubtitle{Ajoitus}
\appear{}

\seminormal
\begin{itemize}
	\item \CDAlert[red]{Arkeologiassa}\appear{, ${ }^{14} \mathrm{C}$-ajoitus on tärkeä työkalu}
	\item ${ }^{14} \mathrm{C}$ on $\beta^{-}$-aktiivinen\appear{, sen puoliintumisaika on $T_{1 / 2}=5730$ vuotta}

\item \CDAlert[blue]{Luonnon hiilenkierron takia}\appear{, suhde ${ }^{14} \mathrm{C} /{ }^{12} \mathrm{C}$ on elävissä organismeissa sama kuin ilmakehässä} \appear{(nykyisin suhde on ${ }^{14} \mathrm{C} /{ }^{12} \mathrm{C}=1,35 \times 10^{-12})$} 
\item Organismi kuolee\appear{, ilmakehän yhdisteet kuten $\mathrm{CO}_{2}$} \appear{tai hiilivedyt} \appear{lopettavat \CDAlert[fuchsia]{hiilenkierron}} 
\implies{ Hiili ei enää uudistu\appear{, ja radioaktiivinen ${ }^{14} \mathrm{C}$ alkaa hajoamaan}}

\item Näytteessä olevan yhden hiiligramman \CDAlert[red]{aktiivisuus} \appear{(vaimenemisnopeus)} \appear{mahdollistaa \CDAlert[red]{kuolinhetken} laskemisen.} \appear{Käytännössä ${ }^{14} \mathrm{C} /{ }^{12} \mathrm{C}$ suhde} \appear{on muuttunut vuosien saatossa}\appear{, heikentäen radiohiiliajoituksen tarkkuutta} \vspace{2mm}

\begin{itemize} \small
	\item ${\sim}$9000 vuotta sitten suhdeluku oli $1,5$-kertainen nykyiseen verrattuna
	\item Ydinkokeet ovat kasvattaneet suhdetta 
	\item \Alert[fuchsia]{Fossiilisten polttoaineiden käyttö on pienentänyt suhdetta} \appear{(merkkejä \CDAlert[fuchsia]{antroposeenistä?})}
\end{itemize}

\end{itemize}

\endofslide 
%\endofsection
\end{frame}
%%%%%%%%%%%%%%%

%%%%%%%%%%%%%%%
\begin{frame}
\frametitle{\texorpdfstring{\culine[\sectiontitleulinecolor]{\sectiontitle}}{\sectiontitle} }
\framesubtitle{Ajoitus}
\appear{}

\begin{itemize}
	\item Arvioidaan seuraavaksi \CDAlert[blue]{maapallon ikä}!

\item Radioaktiivisessa hajoamisessa emoydin $(A)$ \appear{hajoaa reaktiotuotteeksi $(B)$.}\appear{ Näiden ytimien lukumäärät ovat vastaavasti:} 
\[
N_{A}(t)  \equal \appear{N_{0}} \appear{\rme^{-\lambda t}} \qquad \appear{, \qquad N_{B}(t)}  \equal \appear{N_{0}-N_{A}(t)}
\]

\item Tällöin jokaisella ajan hetkellä $t$ on voimassa 
\[
t \equal \frac{1}{\lambda} \appear{\ln} \LP \frac{N_{0}}{N_{A}} \;\RP  \equal \frac{T_{1 / 2}}{\ln 2} \appear{\ln} \LP \frac{N_{0}}{N_{A}} \,\RP  \equal \frac{T_{1 / 2}}{\ln 2} \appear{\ln} \LP \appear{1} \plus \frac{N_{B}}{N_{A}} \,\RP
\]

\item Joidenkin isotooppien pitoisuuksien suhteita \appear{(}$\appear{{ }^{238} \mathrm{U} /{ }^{206} \mathrm{Pb}}\appear{,~{ }^{87} \mathrm{Rb} /{ }^{87} \mathrm{Sr}}\appear{,~{ }^{40} \mathrm{K} /{ }^{40} \mathrm{Ar}}$\appear{)} \appear{käytetään 'kivikelloina'.}\appear{ Niiden avulla on arvioitu \CDAlert[blue]{maanpinnan kivikerroksien iän olevan noin} \CDAlert[blue]{$4,5 \times 10^{9}$ vuotta}}
\end{itemize}

\endofslide 
%\endofsection
\end{frame}
%%%%%%%%%%%%%%%

%%%%%%%%%%%%%%%
\begin{frame}
\frametitle{\texorpdfstring{\culine[\sectiontitleulinecolor]{\sectiontitle}}{\sectiontitle} }
\framesubtitle{Hälyttimet}
\appear{}

% Radioactive nuclei in fire alarms
\begin{itemize}
	\item Am-241 on neptuniumin hajoamisketjuun kuuluva isotooppi
	
	\item Am-241 käytetään perinteisissä kotitalouksien palohälyttimissä ionisoimaan ilmaa

\item \CDAlert[red]{Puhtaassa ilmassa} ilman ionit liikkuvat helposti elektrodien välillä \appear{muodostaen pienen ylläpitovirran}
\implies{ tätä \CDAlert[red]{ylläpitovirtaa mitataan jatkuvasti} }

\item \CDAlert[blue]{Jos savuhiukkasia saapuu} kammioon\appear{, ne poistavat ioneita ilmasta (reagoivat ionisoituneiden kaasumolekyylien kanssa)}\appear{, ilman sähköinen johtavuus laskee} 
\implies{ ylläpitovirta heikkenee/häviää}
\end{itemize}

\endofslide 
%\endofsection
\end{frame}
%%%%%%%%%%%%%%%

%%%%%%%%%%%%%%%
\begin{frame}
\frametitle{\texorpdfstring{\culine[\sectiontitleulinecolor]{\sectiontitle}}{\sectiontitle} }
\appear{}

\begin{itemize}
	\item Hiukkasten synnyttämä \CDAlert[red]{röntgensäteily} ja \CDAlert[blue]{gammasäteily} (PIXE ja PIGE)
\item \CDAlert[red]{PIXE}: 
\item[] \begin{itemize}
	 \item Näytteiden \CDAlert[red]{hivenalkuainepitoisuuksien} alkuaineanalyysi \appear{($Z>20$)} 
	 \item Näytettä pommitetaan matalan energian ioneilla \appear{(${\sim}$MeV)}\appear{, jolloin atomin \CDAlert[red]{K- ja L-kuorien elektronit} emittoivat} \appear{\CDAlert[red]{röntgensäteitä}} 
	 
	 \item Tämän avulla näytteen alkuaineita voidaan tunnistaa 
	 
	 \item Käytettävistä ioneista riippuen \appear{ menetelmä \Alert[fuchsia]{voi soveltua hyvin pintojen tutkimiseen} }
	 
	 \item PIXE:ä käytetään selvittämään puiden alkuainekoostumuksia \appear{(joilla on vuosilustot/-renkaat)} \appear{tai varmistamaan vanhojen maalauksien aitous} \appear{(maaliaineiden koostumus on vaihdellut aikakausittain)}
\end{itemize}

%\item \CDAlert[blue]{PIGE}: Bombarded with higher ion energies, nuclei start emitting $\gamma$-rays from the nucleus in a similar way. Then also lighter elements can be detected (PIGE = Particle-Induced Gamma-ray Emission). For example, the ${ }^{23} \mathrm{Na}\left(\mathrm{p}, \mathrm{p}^{\prime} \mathrm{v}\right)^{23} \mathrm{Na}$ reaction is frequently used for determining bulk sodium content of glasses 
%
%\item PIXE and PIGE are used widely in elemental analysis in material and environmental sciences.
\end{itemize}

\endofslide 
%\endofsection
\end{frame}
%%%%%%%%%%%%%%%

%%%%%%%%%%%%%%%
\begin{frame}
\frametitle{\texorpdfstring{\culine[\sectiontitleulinecolor]{\sectiontitle}}{\sectiontitle} }

\begin{itemize}
	\item Hiukkasten synnyttämä \CDAlert[red]{röntgensäteily} ja \CDAlert[blue]{gammasäteily} (PIXE ja PIGE)

	\item \CDAlert[blue]{PIGE}: 
	\item[] \begin{itemize}
	 \item Kun ytimiä pommitetaan energeettisemmillä ioneilla\appear{, ydin alkaa emittoimaan \Alert[blue]{gammasäteitä}} \appear{samalla tavalla kuin PIXE:ssä}
	 \item Nyt myös kevyempiä alkuaineita voidaan mitata ja tunnistaa 
	  
	 \item Esimerkiksi $\appear{{ }^{23} \mathrm{Na}}\appear{( \mathrm{p}}\appear{,} \appear{\mathrm{p}^{\prime}} \appear{+ \gamma )} \appear{{}^{23} \mathrm{Na}}$ \appear{reaktiota käytetään jatkuvasti \CDAlert[blue]{määrittelemään lasien natriumpitoisuuksia}} 
\end{itemize}

\item[]
\item PIXE:ä ja PIGE:ä käytetään paljon \appear{\CDAlert[fuchsia]{materiaalitutkimuksessa ja ympäristötieteissä}} \appear{\CDAlert[fuchsia]{alkuainekoostumuksien selvittämiseen}} 
\end{itemize}

%\endofslide 
\endofsection
\end{frame}
%%%%%%%%%%%%%%%

